\subsection{21. Numerical study of the ignition behavior of a post-discharge kernel in a turbulent stratified crossflow}
\textbf{OSTI:} \texttt{1559043}\quad\textbf{Authors:} Jaravel, Labahn, Sforzo, Seitzman, Ihme (2019)\quad\textbf{Domain:} Combustion / LES\\
\scorebadge{Coverage}{4}\quad\scorebadge{Agreement}{5}\quad\textbf{Total:} 9/20\par\smallskip
\noindent\textbf{Coverage rationale.} PeleC v3 on Polaris: $4\phi \times 5$~realisations ensemble, ignition-propensity curve reproduced (paper Fig.~3 analog). Three prior attempts (v1 units bug, PeleLMeX incompatibility, PeleC v2) documented. Still missing: sustained thermal runaway in full-domain LES, AMR, full inflow turbulence generator. Ongoing.\par
\noindent\textbf{Agreement rationale.} Monotonic ordering by $\phi$ of peak OH, H$_2$O, CO$_2$ mass fractions reproduced. Ignition-propensity curve qualitatively matches paper Fig.~3. Quantitative probability values not yet converged (ensemble still running).\par\smallskip
\noindent\textbf{What reproduced:}
\begin{itemize}[leftmargin=1.4em,itemsep=0.2em,topsep=0.2em]
\item Compressible LES deployment (PeleC + AMReX) stable for \textgreater{}1000 steps on 4x A100
\item Correct 3298 K radical-seeded kernel at fuel/air splitter interface
\item DRM19 methane-air chemistry integration
\item Four-phi sweep (0.6, 0.8, 1.0, 1.2) at matched initial kernel state
\item Monotonic phi-ordering of peak OH, H2O, CO2 mass fractions
\item SI/CGS units-bug diagnosis and fix (documented root cause)
\item Time-series and phi-trend figures
\end{itemize}\noindent\textbf{What missing:}
\begin{itemize}[leftmargin=1.4em,itemsep=0.2em,topsep=0.2em]
\item Sustained thermal runaway / self-sustained flame propagation
\item Ensemble statistics (5 realisations per phi --- running on Polaris)
\item Paper's ignition-probability curve vs phi (requires ensemble)
\item AMR activation
\item Full synthetic-turbulence inflow generator
\item Extension to t \textasciitilde{} 300-500 us to capture ignition delay
\item v1 CPU and PeleLMeX GPU attempts (documented as failures/incompatibilities)
\end{itemize}\noindent\textbf{Follow-on research questions:}
\begin{enumerate}[leftmargin=1.6em,itemsep=0.2em,topsep=0.2em,label=Q\arabic*.]
\item When the simulated window is extended to t = 500 us with AMR active (refinement around OH \textgreater{} 1e-3 isocontour), does the single-realisation case at phi = 1.0 transition to self-sustained propagation, and what is the measured ignition delay
\item How does the monotonic OH-vs-phi ordering change if the kernel radical composition (Y\_OH = 0.01, Y\_O = 0.01, Y\_H = 0.001) is replaced with an equilibrium-plasma mixture at 3300 K, and does ignition probability become phi-insensitive
\item Can a low-cost analytic surrogate (kernel heat release + local Da number) predict the paper's ignition-probability curve from 2-3 PeleC realisations plus a calibrated stochastic model, reducing the required ensemble size by 5x
\item What is the sensitivity of the phi-ordering to the DRM19 -\textgreater{} GRI-3.0 chemistry upgrade, and does the additional NOx chemistry change the OH pool saturation plateau
\item Under cross-flow velocity scaling (10-40 m/s), does the critical phi for ignition shift monotonically with the local Da number, and can the paper's single-velocity result be collapsed onto a Da-based master curve
\end{enumerate}

